% ----------------------------------------------------------------------
% Proposition: neural-network approximation on a finite probability space
% This file is an *artifact* cited in Section 4.2 of the thesis. It reproduces
% Buehler et al. (2019), Proposition 4.3, and tags *exactly* the three
% places where finiteness of Omega is used (F1) to (F3).
% When integrating into the final thesis, strip the standalone preamble
% and include this under the same numbering as the surrounding document
% (keep the "% draft: ..." comment to preserve the original label).
% Cross-references below are faked (hardcoded) so this artifact compiles
% without defining the targets; comments say what they should point to.
% ----------------------------------------------------------------------
\documentclass[11pt,a4paper]{report}
\usepackage[T1]{fontenc}
\usepackage[utf8]{inputenc}
\usepackage{amsmath,amssymb,amsthm}
\usepackage{enumitem}

\numberwithin{equation}{chapter}

\theoremstyle{plain}
\newtheorem{proposition}{Proposition}[chapter]
\newtheorem{lemma}[proposition]{Lemma}
\newtheorem{theorem}[proposition]{Theorem}
\newtheorem{corollary}[proposition]{Corollary}
\theoremstyle{definition}
\newtheorem{assumption}[proposition]{Assumption}
\newtheorem{definition}[proposition]{Definition}
\newtheorem{example}[proposition]{Example}
% Upright body, bold head (not the italic amsthm "remark" style).
\newtheorem{remark}[proposition]{Remark}

\newcommand{\E}{\mathbb{E}}
\newcommand{\PP}{\mathbb{P}}
\newcommand{\R}{\mathbb{R}}
\newcommand{\N}{\mathbb{N}}
\newcommand{\F}{\mathcal{F}}
\newcommand{\HH}{\mathcal{H}}
\newcommand{\Hu}{\mathcal{H}^{\mathrm{u}}}
\newcommand{\NNet}{\mathcal{N\!N}}
\newcommand{\as}{\ \PP\text{-a.s.}}

\begin{document}

\setcounter{chapter}{3}

\chapter{Artifact: finite-space approximation}

\subsection*{The approximation theorem on a finite probability space}

% draft: Proposition 4.3 (Buehler et al.)
% Attribution in prose (styling.md, item 21): no optional title on the environment.
B\"uhler et al., Proposition~4.3, prove the following on a finite probability space. In the proof we flag each use of finiteness of $\Omega$ as (F1) to (F3); these are exactly the points that need attention on a general probability space.
\begin{proposition}\label{prop:finitespace}
  Suppose the probability space $\Omega=\{\omega_1,\dots,\omega_N\}$ is
  finite with $\PP[\{\omega_i\}]>0$ for all $i$, and
  $\rho:\R^N\to\R$ is a convex risk measure. Then for every
  claim $X$ one has $\displaystyle\lim_{M\to\infty}\pi_M(X)=\pi(X)$.
\end{proposition}

\begin{proof}
  % FAKE REF: "Lemma~4.1" -> replace by \ref{lem:monotone}
  %   target: monotone sandwich lemma (ch4_structure.md, Lemma 4.1 / end of §4.1;
  %   same statement as lem:monotone in extension_of_approximation_on_general.tex,
  %   artifact lemma4_1_monotone.tex).
  By monotonicity of $(\pi_M)$ (Lemma~4.1) it suffices to
  show that for every $\varepsilon>0$ some $M$ satisfies
  $\pi_M(X)\le\pi(X)+\varepsilon$.

  Choose $\delta\in\Hu$ with
  % label kept for merge; not used by \ref/\eqref in this artifact
  \begin{equation}\label{eq:finitespace:nearopt}
    \rho\bigl(X+W(H\circ\delta)\bigr)\le\pi(X)+\frac{\varepsilon}{2}.
  \end{equation}
  Since each $\delta_k$ is $\F_k$-measurable with
  $\F_k=\sigma(I_0,\dots,I_k)$, and since $\delta_{k-1}$ is a
  function of $(I_0,\dots,I_{k-1}, \delta_{k-2})$ we can recusively define $f_k$ so that
  $\delta_k=f_k(I_0,\dots,I_k)$.

  % Run-in italic labels (not \paragraph): ordinary paragraphs => indent, no extra skip.
  % \indent: amsthm proof often suppresses the first-line indent.
  \indent\textit{(F1) Integrability of the hedge ratios.}
  Because $\Omega$ is finite, each $\delta_k$ takes finitely many values
  and every component of $f_k$ lies in $L^1(\mu_k)$ where $\mu_k$ is the law of
  $(I_0,\dots,I_k)$. Finally, Hornik's theorem yields networks
  $F_{k,m}\to f_k$ in $L^1(\mu_k)$.

  \indent\textit{(F2) From $L^1$ to pointwise convergence.}
  On a finite $\Omega$, $L^1$ convergence implies pointwise convergence:
  $\delta^m_k:=F_{k,m}(I_0,\dots,I_k)\to\delta_k$ for every
  $\omega\in\Omega$.

  \indent Set $G_m:=((H\circ\delta^m)\cdot S)_T$ and $C_m:=C_T(H\circ\delta^m)$.
  Since $H$ and the stochastic integral are continuous, $G_m\to G$ pointwise.
  Furthermore, using upper semicontinuity of the $c_k$ and continuity of $H$ again,
  one has $\limsup_m C_m=\overline C\le C:=C_T(H\circ\delta)$ pointwise.

  \indent\textit{(F3) Continuity of the risk measure.}
  Identifying random variables with vectors in $\R^N$, the convex
  map $\rho:\R^N\to\R$ is continuous. Continuity and
  monotonicity of $\rho$ give
  \[
    \lim_{m\to\infty}\rho(X+G_m-C_m)=\rho(X+G-\overline C)
    \le\rho(X+G-C)\le\pi(X)+\frac{\varepsilon}{2}
  \]
  % FAKE REF: "(4.1)" -> replace by \eqref{eq:finitespace:nearopt}
  %   target: the near-optimal inequality displayed just above in this proof.
  by~(4.1). Pick $m$ large enough so that the
  left-hand side is $\le\pi(X)+\varepsilon$. Because $\delta^m$ belongs
  to $\HH_M$ for all $M$ once $m$ is large enough (nested network
  classes), $\pi_M(X)\le\pi(X)+\varepsilon$ follows.
\end{proof}

\subsection*{What breaks on a general probability space}

\begin{remark}
  \textbf{(F1): hedge ratios need not be integrable.}
  On general $\Omega$, the representing functions $f_k$ need not lie in
  $L^1(\mu_k)$, and then Hornik's theorem cannot even be invoked: $f_k$ is
  not an element of the space in which networks are dense.
\end{remark}

\begin{example}\label{ex:integrability}
  Let $I_0$ have a standard Cauchy law $\mu_0$ and $\F_0=\sigma(I_0)$.
  The strategy $\delta_0=f(I_0)$ with $f(x)=|x|$ is perfectly admissible
  as a measurable adapted process, but $f\notin L^1(\mu_0)$, so no
  sequence of $\mu_0$-integrable functions (in particular no sequence
  of networks with bounded activation, which are bounded functions)
  converges to $f$ in $L^1(\mu_0)$.
\end{example}

\begin{remark}
  \textbf{(F2): only subsequential almost sure convergence (harmless).}
  On general $\Omega$, $L^1(\PP)$-convergence yields almost sure
  convergence only along a subsequence. This is harmless for the
  argument: the proof only requires \emph{one} good approximating
  strategy, so we may freely pass to subsequences.
\end{remark}

\begin{remark}
  \textbf{(F3): the risk measure is not continuous along almost sure
    convergence.}
  This is the decisive failure. On $\R^N$, convexity of $\rho$ implies
  continuity. On an infinite-dimensional domain, a convex risk measure
  need not be continuous along almost sure convergence. The direction needed
  in the final step of the proof is an \emph{upper} semicontinuity:
  \[
    \limsup_{m\to\infty}\rho(Y_m)\le\rho(Y)\qquad\text{whenever }Y_m\to Y\as
  \]
  % FAKE REF: "(A3)" -> the half-bounded Fatou property of Chapter 3,
  %   Section 3.4; lower semicontinuity, i.e. the opposite direction.
  This fails even for the entropic risk measure, which satisfies the
  half-bounded \emph{lower} semicontinuity (A3) of Chapter~3.
\end{remark}

\begin{example}\label{ex:explosion}
  Work on $\Omega=(0,1)$ with the Borel $\sigma$-field and Lebesgue
  measure, let $\lambda>0$ and let
  $\rho_\lambda(Y)=\frac{1}{\lambda}\log\E[e^{-\lambda Y}]$ be the
  entropic risk measure. Set $A_m:=(0,e^{-m})$ and
  \[
    Y_m:=-\frac{2m}{\lambda}\mathbf{1}_{A_m}.
  \]
  For every $\omega\in(0,1)$ one has $\omega\notin A_m$ for all
  $m>\log(1/\omega)$, so $Y_m\to0$ everywhere. But
  \[
    \E\bigl[e^{-\lambda Y_m}\bigr]=1-e^{-m}+e^{-m}e^{2m}
    =1-e^{-m}+e^{m}\longrightarrow\infty,
  \]
  so $\rho_\lambda(Y_m)\to\infty$ while $\rho_\lambda(0)=0$. Small
  events carrying large losses make the risk explode along an almost
  surely vanishing sequence.
\end{example}

\end{document}
